% manuscritos/plan-de-tesis/main.tex -- Plan de tesis de Maestria.
%
% Compilar desde PowerShell (no Git Bash: MSYS reescribe BIBINPUTS y bibtex no encuentra el .bib):
%   .\manuscritos\herramientas\compilar.ps1 -Documento plan-de-tesis
% Salida: manuscritos/plan-de-tesis/build/main.pdf (fuera de git).
%
% ESTRUCTURA: calcada del plan de tesis aprobado que el autor aporto como ejemplo,
%   manuscritos/info-institucional/plan-tesis/Ejemplos/Plan Tesis de Maestria - Joaquin Ignacio
%   Aramendia_v2.pdf (4 capitulos, ~10 paginas de cuerpo + bibliografia).
% Requisitos de forma: manuscritos/info-institucional/plan-tesis/Instructivo _Aprobacion plan
%   tesis + Director (estudiantes).docx y 2-Caratula-PLAN de tesis.docx.
\documentclass{../comun/tesisuba}

\etapa{Plan de Tesis}
\facultadsecundaria{Facultad de Ingenier\'ia}

% El instructivo pide el titulo sin mayusculas salvo nombres propios, y avisa que ya en esta
% instancia se mira con atencion porque arrastra a las siguientes. PROVISORIO hasta que el autor
% y el director lo cierren.
\titulo{Pron\'ostico de caudal\\ del r\'io Uruguay\\ con funciones\\ de p\'erdida asim\'etricas}
\leyenda{Plan de tesis presentado para optar al t\'itulo de Mag\'ister\\
en Explotaci\'on de Datos y Descubrimiento de Conocimiento}
\autor{Tesista: Esp. Joaquín Sebastian Tschopp}
\director{Gustavo Denicolay}
% FALTA: el autor confirma nombre y grado academico del codirector.
\codirector{}
\lugardetrabajo{}
% FALTA: fecha de elevacion de la nota al director de la maestria.
\fecha{Buenos Aires, \ldots\ de 2026}

\begin{document}

\portada

\tableofcontents

\newpage

% Plan de tesis -- Bloque administrativo, previo a los capitulos.
%
% Lo exige el punto 2.b del instructivo de aprobacion del plan: "Dentro del plan detallar los datos
% del Director y Codirector (si corresponde), aclarando su grado academico y filiacion
% institucional actual, junto con los correos electronicos de contacto".
% El plan de ejemplo no trae este bloque; se agrega porque el instructivo lo pide de forma expresa.
\section*{Dirección de la tesis}
\addcontentsline{toc}{section}{Dirección de la tesis}

\noindent
\begin{tabular}{@{}lp{8.5cm}@{}}
  \toprule
  \multicolumn{2}{@{}l}{\textbf{Tesista}} \\
  Nombre y apellido      & Joaquín Sebastian Tschopp \\
  Grado académico        & Especialista en Explotación de Datos y Descubrimiento de Conocimiento \\
  Filiación institucional & \textit{(falta: completar)} \\
  Correo electrónico     & \textit{(falta: completar)} \\
  \midrule
  \multicolumn{2}{@{}l}{\textbf{Director}} \\
  Nombre y apellido      & Gustavo Denicolay Pacheco \\
  Grado académico        & \textit{(falta: completar)} \\
  Filiación institucional & \textit{(falta: completar)} \\
  Correo electrónico     & \textit{(falta: completar)} \\
  \midrule
  \multicolumn{2}{@{}l}{\textbf{Codirector}} \\
  Nombre y apellido      & \textit{(falta: completar)} \\
  Grado académico        & \textit{(falta: completar)} \\
  Filiación institucional & \textit{(falta: completar)} \\
  Correo electrónico     & \textit{(falta: completar)} \\
  \bottomrule
\end{tabular}

\newpage

% Plan de tesis -- Capitulo 1.
%
% Correlato en el plan de ejemplo: capitulo 1 (Introduccion), con 1.1 Descripcion del problema y
% motivacion, 1.2 Trabajos previos y 1.3 Objetivos. Extension del ejemplo: ~3 paginas.
\chapter{Introducción}

\section{Descripción del problema y motivación}

% PROVISORIO. Esbozo para discutir con el director; las afirmaciones cuantitativas todavia no
% estan referenciadas y hay que cerrarlas contra fuente antes de elevar el plan.
El pronóstico de caudal a pocos días de anticipación es la información que ordena la operación de
un aprovechamiento hidroeléctrico: define cuánta agua se turbina, cuánta se eroga y con cuánta
antelación se avisa aguas abajo.
En el río Uruguay esa decisión se toma sobre una cuenca de respuesta rápida, donde el tiempo entre
la lluvia en las nacientes y el paso de la onda de crecida por el embalse es del orden de los días.

El problema no es solamente predecir bien en promedio.
Los errores no cuestan lo mismo en las dos direcciones ni en todos los regímenes: subestimar una
crecida obliga a erogar de apuro y traslada el problema aguas abajo, mientras que sobreestimarla
tiene un costo de oportunidad en generación.
Sin embargo, los modelos de pronóstico de caudal se entrenan casi siempre minimizando un error
simétrico --- error cuadrático medio, o su versión normalizada como el coeficiente de
Nash--Sutcliffe --- que trata ambos errores como equivalentes.

Esta tesis parte de esa brecha: la función que se optimiza durante el entrenamiento no representa
el costo real de la decisión que el pronóstico alimenta.

% FALTA: cerrar la motivacion con el encuadre institucional (el autor trabaja en un aprovechamiento
% de la cuenca) una vez que el autor defina cuanto de eso quiere que aparezca en el documento.

\section{Trabajos previos}

% BORRADOR PROVISORIO del 2026-09-18, escrito sobre el relevamiento de 71 referencias que vive en
% research/catalog/. Los metadatos de todas las citas estan verificados contra Crossref, pero las
% AFIRMACIONES sobre el contenido de cada trabajo salen del relevamiento y no de la lectura del
% paper: hay que confirmarlas contra la fuente antes de elevar el plan. Ver PLAN.md, seccion 5.

\subsection{Modelos en operación}

El pronóstico hidrológico operativo se reparte hoy entre dos linajes.
El primero es físico y distribuido, con dos décadas de rodaje: GloFAS acopla los ensambles
meteorológicos del Centro Europeo con un modelo hidrológico distribuido para dar alerta de
crecidas a escala global \cite{alfieri-2013-glofas}; el \textit{National Water Model} de la NOAA
hace lo propio sobre los Estados Unidos y documenta además cómo un sistema en producción trata
los embalses y la regulación antrópica \cite{cosgrove-2024-national-water-model}; y sobre esta
región el modelo continental MGB cubre explícitamente la cuenca del Plata
\cite{siqueira-2018-mgb-sudamerica}.
El aporte más transferible de esta familia no es el modelo sino el protocolo de verificación:
reportar el desempeño contra líneas de base de persistencia y de climatología, y mostrar cómo se
degrada con el horizonte \cite{harrigan-2023-glofas-reforecast}.

El segundo linaje es aprendizaje profundo, y llegó a producción muy rápido.
La red LSTM aplicada a lluvia-escorrentía iguala o supera a los modelos conceptuales calibrados
cuenca por cuenca \cite{kratzert-2018-lstm-lluvia-escorrentia}, y entrenada sobre muchas cuencas a
la vez aprende comportamiento compartido que luego modula con atributos locales
\cite{kratzert-2019-universal-regional-local}.
Sobre esa base se construyó un sistema de alerta de crecidas que hoy opera en decenas de países
\cite{nevo-2022-pronostico-operativo-ml, nearing-2024-global-extreme-floods}.
Importa registrar que la evaluación de ese sistema está en disputa: se le objetó que las
decisiones metodológicas de su evaluación --- cómo se define un evento extremo, qué ventana se
usa, contra qué se compara --- inflaron las métricas reportadas
\cite{li-2026-critica-despliegue-ia}.
La lección para este trabajo no es tomar partido, sino declarar y defender cada decisión de
evaluación en lugar de darla por obvia.

Sobre las arquitecturas más nuevas, la evidencia es menos entusiasta de lo que suele suponerse.
En series temporales generales, un modelo lineal de una sola capa iguala o supera a toda la
familia de transformadores de horizonte largo \cite{zeng-2023-dlinear}.
En hidrología, el transformador común no resulta competitivo frente a la LSTM, y falla
especialmente en caudales altos \cite{liu-2024-transformer-limite-hidrologico}.
Los modelos fundacionales de series temporales, evaluados sin entrenamiento específico, todavía
quedan por debajo de una LSTM entrenada \cite{sun-2026-tsfm-zeroshot-caudal}.
El matiz importante es que la ventaja de la LSTM se concentra en los horizontes cortos, y los
mecanismos de atención ganan terreno a medida que el horizonte se alarga
\cite{liu-2025-rnn-a-transformers}; combinada con recurrencia, la atención sí mejora el
pronóstico de caudal \cite{rasiyakoya-2024-tft-caudal}, y existe una arquitectura pensada
exactamente para este problema, multi-horizonte y con covariables conocidas a futuro
\cite{lim-2021-tft}, que es la situación que crea disponer del pronóstico de lluvia.
La ventana de 1 a 7 días de este trabajo cae justo en esa zona de transición.

Para datos tabulares, que es la forma que tiene el conjunto de datos de esta tesis, los modelos de
árboles siguen siendo estado del arte frente a las redes profundas
\cite{grinsztajn-2022-arboles-tabular}, dominaron la competencia M5
\cite{makridakis-2022-m5-accuracy} y ganaron una competencia hidrológica con evaluación externa
\cite{drivendata-2025-water-supply-rodeo}.
La contracara también está documentada: con más datos la LSTM los supera
\cite{gauch-2021-camels-datos-limitados}, y los árboles no extrapolan fuera del rango observado
\cite{januschowski-2022-forecasting-trees}, que es precisamente el régimen de las crecidas
extremas.

Dos advertencias metodológicas quedan pendientes de respuesta.
La primera es que entrenar una LSTM sobre una sola cuenca está explícitamente desaconsejado
\cite{kratzert-2024-nunca-una-sola-cuenca}; la salida publicada es preentrenar sobre un conjunto
grande y ajustar localmente \cite{ryd-2025-fine-tuning-local}, para lo cual existen los datos
\cite{chagas-2020-camels-br, kratzert-2023-caravan}.
La segunda es que entrenar con lluvia observada y operar con lluvia pronosticada abre una brecha
de desempeño real, para la que también hay una receta publicada \cite{taccari-2026-aifl}.

Sobre el área de estudio, el antecedente más cercano evalúa un sistema operativo de pronóstico de
caudal en el alto río Uruguay, corrido a diario con diez días de anticipación, y reporta que la
ocurrencia de crecidas se anticipa con dos o tres días útiles
\cite{fan-2017-uruguai-operacional}; la verificación de afluencias a embalses de la región con
ensambles meteorológicos tiene tratamiento propio \cite{fan-2015-tigge-embalses}.
En cuanto a aprendizaje automático sobre el propio río, lo publicado predice \emph{nivel} a tres
días \cite{mattiuzi-2021-m5-uruguai}, y el producto operativo vigente en la cuenca es también de
nivel y de muy corto plazo \cite{mattiuzi-2023-sah-uruguai}.
Para las afluencias a Salto Grande, el trabajo disponible opera en escala mensual y estacional
\cite{talento-2013-salto-grande}.
% AFIRMACION FUERTE, a sostener o retirar: de este relevamiento no surge ningun trabajo publicado
% de pronostico de CAUDAL del rio Uruguay con aprendizaje automatico en horizontes de 1 a 7 dias.
% Es el hueco que justifica la tesis. Conviene que el director y el codirector lo confirmen antes
% de escribirlo en el plan.

\subsection{Funciones objetivo}

La teoría moderna de evaluación de pronósticos fija una regla que ordena este campo: una función
de pérdida sirve para comparar modelos sólo si es \textit{consistente} para algún funcional
estadístico, y ese funcional debe declararse de antemano
\cite{gneiting-2011-making-evaluating-point-forecasts, gneiting-raftery-2007-scoring-rules}.
El error cuadrático es consistente para la media; la pérdida \textit{pinball}, para el cuantil
\cite{koenker-bassett-1978}; la cuadrática asimétrica, para el expectil
\cite{newey-powell-1987}; y la de Huber asimétrica, para un funcional intermedio entre los dos
\cite{taggart-2022-generalized-huber-loss}.
Toda pérdida consistente para un cuantil o un expectil se descompone en \textit{scores}
elementales, lo que habilita comparaciones robustas \cite{ehm-2016-quantiles-expectiles}.
De ahí se desprende la principal dificultad metodológica de este trabajo: cambiar la asimetría
cambia el funcional que se estima, y no sólo el peso de los errores, de modo que comparar con una
métrica única a modelos entrenados con pérdidas distintas es evaluarlos con la directiva
equivocada.

En hidrología las funciones objetivo tienen su propia genealogía.
El coeficiente de Nash--Sutcliffe \cite{nash-sutcliffe-1970} es el error cuadrático reescalado por
la varianza observada, y por lo tanto estima el mismo funcional.
El de Kling--Gupta lo descompone en correlación, sesgo y variabilidad \cite{gupta-2009-kge}, y
tiene variantes sucesivas que corrigen sus supuestos \cite{kling-2012-modified-kge,
pool-2018-nonparametric-kge}; no es, en cambio, una pérdida consistente para ningún funcional
conocido.
Ambos índices tienen una incertidumbre muestral grande, porque unos pocos días de crecida dominan
su valor \cite{clark-2021-abuse-performance-metrics}: es un argumento contra ordenar modelos por
un único número.
Las transformaciones del caudal son, de hecho, un cambio encubierto de función de pérdida
\cite{santos-2018-log-kge-pitfalls, thirel-2024-streamflow-transformations}, algo que recién ahora
se formalizó en el lenguaje de la consistencia \cite{tyralis-2026-variable-transformations-loss};
y volver de la escala logarítmica a metros cúbicos por segundo introduce un sesgo de
retransformación conocido \cite{duan-1983-smearing-retransformation}.

La familia asimétrica tiene respaldo teórico y algunos antecedentes empíricos.
El resultado clásico es que bajo pérdida asimétrica el predictor óptimo es sesgado, y su sesgo
varía en el tiempo \cite{christoffersen-diebold-1997-optimal-prediction-asymmetric}: el sesgo es
el objetivo, no un defecto.
La alternativa suave es la pérdida LINEX \cite{zellner-1986-asymmetric-loss-linex}.
En hidrología ya se usaron funciones de error asimétricas para fijar umbrales de alerta
\cite{toth-2016-asymmetric-flood-warning}, se propuso una pérdida asimétrica de picos
\cite{tofighi-2025-asymmetric-peak-loss}, y se entrenaron redes con ocho funciones de pérdida
distintas, incluidas expectílicas
\cite{dahal-2026-ensemble-diverse-loss-functions}.
La regresión expectílica en hidrología ya tiene tratamiento propio
\cite{tyralis-2023-expectile-hydrological} y una implementación en redes profundas que anida los
casos cuantílico y expectílico \cite{tyralis-2025-deep-huber-quantile-networks}; si se opta por
estimar varios cuantiles a la vez, hay que resolver que no se crucen
\cite{cannon-2018-mcqrnn-noncrossing}.
Conviene además registrar que los sistemas operativos más recientes no usan una pérdida simétrica
sino distribuciones asimétricas aprendidas \cite{klotz-2022-uncertainty-deep-learning-rainfall-runoff,
nearing-2024-global-extreme-floods}: la asimetría ya está en producción, y el aporte de esta tesis
tiene que ser más específico que introducirla.
Fuera de la hidrología, hay evidencia fuerte en energía de que una pérdida asimétrica calibrada
con el cociente de costos real reduce el costo de operación
\cite{wu-2021-asymmetric-svr-load-forecasting}, y de que la pérdida que maximiza el beneficio no es
la que minimiza el error cuadrático \cite{serafin-weron-2025-loss-functions-trading}.

El eslabón que conecta la función de pérdida con la operación es la teoría de la decisión.
La relación costo-pérdida establece que el nivel de cuantil óptimo es el cociente entre el costo
de actuar y la pérdida de no hacerlo \cite{richardson-2000-relative-economic-value}, criterio ya
aplicado a sistemas reales de alerta \cite{verkade-werner-2011-flood-warning-benefits} y discutido
para decisiones hidrológicas con aversión al riesgo \cite{matte-2017-beyond-cost-loss}.
Dicho de otro modo, el nivel de asimetría no es un hiperparámetro: es un dato del problema de
decisión.
La línea que lleva ese razonamiento hasta el final entrena minimizando directamente el costo de la
decisión \cite{elmachtoub-grigas-2022-smart-predict-optimize}.
Hay que registrar, sin embargo, que más precisión no implica más valor en la operación de un
embalse: la relación no es monótona y depende del objetivo operativo
\cite{turner-2017-forecast-skill-value-reservoir}.

Por último, dos límites que este trabajo debe asumir de entrada.
El primero es que evaluar sólo sobre los eventos extremos distorsiona la comparación
\cite{lerch-2017-forecasters-dilemma}; la forma correcta de mirar el desempeño en crecidas es
descomponer un \textit{score} consistente por regiones del rango
\cite{taggart-2022-extreme-events-scoring} o usar \textit{scores} ponderados por umbral que
conserven la propiedad \cite{gneiting-ranjan-2011-weighted-scoring,
allen-2023-transformed-kernel-scores}.
El segundo es un resultado negativo duro: las propiedades de cola no son elicitables
\cite{brehmer-strokorb-2019-tail-properties}.
Una pérdida asimétrica puede mejorar el desempeño en la región alta del rango; no puede prometer
que el modelo aprenda la cola.

\section{Objetivos}

% PROVISORIO: el objetivo general y los especificos los cierra el autor con el director. Lo que
% sigue es la formulacion que se desprende del trabajo ya hecho en el repositorio.

\subsection{Objetivo general}

Evaluar si reemplazar la función de pérdida simétrica habitual por una función asimétrica y
sensible al régimen hidrológico mejora el pronóstico operativo de caudal diario del río Uruguay en
horizontes de 1 a 7 días, respecto de las alternativas documentadas en la literatura y de una
línea de base de persistencia.

\subsection{Objetivos específicos}

\begin{enumerate}[nosep]
  \item Relevar el estado del arte en modelos de pronóstico de caudal efectivamente probados en
        operación, y en funciones de ganancia y pérdida documentadas para este problema.
  \item Construir un conjunto de datos diario reproducible para la cuenca, integrando registros de
        caudal y nivel, precipitación observada y precipitación pronosticada.
  \item Implementar y comparar, bajo un mismo protocolo de validación, las funciones objetivo
        relevadas frente a la función asimétrica propuesta.
  \item Establecer si las diferencias observadas entre modelos son estadísticamente distinguibles,
        y no atribuibles al azar del período de evaluación.
  \item Cuantificar la mejora en términos de la decisión operativa, y no sólo en términos de error
        estadístico.
\end{enumerate}

% Plan de tesis -- Capitulo 2.
%
% Correlato en el plan de ejemplo: capitulo 2 (Materiales y Metodos), con una seccion de base
% tecnica sobre la que se apoya la propuesta, una de Metodos (que es donde el ejemplo enuncia su
% hipotesis) y una de Datos. Extension del ejemplo: ~4 paginas.
\chapter{Materiales y Métodos}

\section{Funciones de pérdida asimétricas}

% Base tecnica sobre la que se apoya la propuesta. Las dos familias y sus referencias canonicas ya
% estan en el catalogo (koenker-bassett-1978, newey-powell-1987, ehm-2016-quantiles-expectiles).
Sea $u = y - \hat{y}$ el error de pronóstico y $\tau \in (0,1)$ un parámetro de asimetría.
La regresión cuantílica \cite{koenker-bassett-1978} minimiza la pérdida \textit{pinball}
\begin{equation}
  \rho^{\mathrm{q}}_{\tau}(u) = u \left( \tau - \mathbf{1}\{u < 0\} \right),
\end{equation}
cuyo minimizador poblacional es el cuantil de nivel $\tau$ de la distribución de $y$.
La regresión expectílica \cite{newey-powell-1987} reemplaza el valor absoluto implícito por un
término cuadrático,
\begin{equation}
  \rho^{\mathrm{e}}_{\tau}(u) = \left| \tau - \mathbf{1}\{u < 0\} \right| u^{2},
\end{equation}
y su minimizador es el expectil de nivel $\tau$.
Para $\tau = 1/2$ ambas se reducen a los casos simétricos conocidos: el error absoluto y el error
cuadrático.
Con $\tau > 1/2$ la pérdida penaliza más la subestimación que la sobreestimación.

La diferencia práctica entre las dos familias es la ponderación de los errores grandes.
El expectil conserva el término cuadrático, de modo que el gradiente crece con la magnitud del
error; eso lo hace derivable en todo el dominio y directamente utilizable como objetivo de
entrenamiento por descenso de gradiente, a diferencia de la pérdida \textit{pinball}, que no es
derivable en el origen.

% PENDIENTE: cerrar esta seccion con la discusion de elicitabilidad --- que funcional estima cada
% perdida y por que eso condiciona la comparacion entre modelos --- cuando este consolidado el
% relevamiento (eje b).

\section{Métodos}

% En el plan de ejemplo esta es la seccion donde se enuncia la hipotesis del trabajo. Se replica.

% PROVISORIO: la hipotesis la cierra el autor con el director.
La hipótesis del trabajo es que el nivel de asimetría conveniente no es constante: en estiaje el
costo de sobreestimar y el de subestimar son comparables, mientras que en crecida la asimetría es
marcada.
Se propone entonces modular el parámetro de asimetría con el régimen hidrológico, es decir,
reemplazar el $\tau$ fijo por
\begin{equation}
  \tau_t = g(\mathbf{x}_t),
\end{equation}
donde $\mathbf{x}_t$ resume el estado de la cuenca en el instante $t$ y $g$ es un modulador
acotado en $(0,1)$.

% PENDIENTE: la forma funcional concreta de g, el conjunto de modelos a comparar y el protocolo de
% validacion se completan con lo que devuelva el relevamiento. Lo ya implementado en el repositorio
% (validacion walk-forward, linea de base de persistencia, contraste de Diebold-Mariano) se
% describe en el capitulo 3.

\section{Datos}

% Inventario verificado contra la documentacion del repositorio (docs/roadmap.md,
% docs/data_sources.md, docs/gold_consolidation_contract.md, docs/gold_quality_report.md).
% Las cifras corresponden a la medicion del 2026-08-26 y hay que refrescarlas antes de elevar.
El conjunto de datos ya está construido y es el insumo de esta tesis.
Cubre la cuenca alta del río Uruguay, aguas arriba de la frontera entre Brasil y Argentina, y se
consolida con una arquitectura por capas sobre \textit{Databricks}.

La variable objetivo es el caudal diario, en m\textsuperscript{3}/s, en la estación 74100000 (Irai)
de la Agência Nacional de Águas de Brasil, sobre la frontera.
El caudal no se mide directamente: se deriva del nivel observado mediante la curva de aforo de la
estación, una ley de potencia $Q = A\,(H - H_0)^{N}$ cuyos coeficientes cambian por tramo y por
período de vigencia.
El nivel se conserva como objetivo secundario, de modo que la serie medida nunca se pierde detrás
de la conversión.
Los horizontes de pronóstico son t+1 a t+7, más t+14.

\begin{table}[H]
  \centering
  \caption{Fuentes de datos del conjunto consolidado.}
  \label{tab:fuentes}
  \small
  \begin{tabular}{@{}llll@{}}
    \toprule
    \textbf{Fuente} & \textbf{Variable} & \textbf{Tipo} & \textbf{Cobertura} \\
    \midrule
    ANA (Brasil)        & Nivel y caudal      & Observado   & desde 1950 \\
    ANA (Brasil)        & Precipitación       & Observado   & desde 1923 \\
    ANA (Brasil)        & Curvas de aforo     & Observado   & desde 1948 \\
    CPTEC/INPE MERGE    & Precipitación       & Observado, grilla 0,1\textdegree{} & desde 1998 \\
    CPTEC/INPE SAMeT    & Temperatura         & Observado, grilla 0,05\textdegree{} & desde 2000 \\
    INMET (Brasil)      & Temperatura         & Observado   & desde 2000 \\
    GEFS Reforecast v12 & Precipitación       & Pronosticado & 2000--2019 \\
    ECMWF TIGGE         & Precipitación       & Pronosticado & desde 2006 \\
    \bottomrule
  \end{tabular}
\end{table}

La distinción entre las dos últimas filas y el resto es la que gobierna el diseño del conjunto: las
variables pronosticadas son las únicas que estarán disponibles en el momento en que el modelo deba
operar, y por eso el histórico de pronósticos --- y no el de observaciones --- es el que fija el
piso temporal del conjunto en el año 2000.
El resultado es una tabla de 9.732 días entre el 1 de enero de 2000 y el 23 de agosto de 2026, con
83 columnas, sin huecos de calendario.
El agregado de precipitación y caudal de la cuenca se construye sobre 36 estaciones.

% FALTA: refrescar las cifras contra Databricks antes de elevar el plan, y decidir con el director
% si conviene declarar las limitaciones de calidad ya documentadas (curvas de aforo extrapoladas en
% algunas estaciones del agregado, desfasaje de la ventana diaria de MERGE contra la de ANA).

% Plan de tesis -- Capitulo 3.
%
% Correlato en el plan de ejemplo: capitulo 3 (Experimentos preliminares y resultados), ~2 paginas.
% El ejemplo muestra una tabla chica con la comparacion contra el metodo base y enumera los
% experimentos siguientes. Es el capitulo que convierte la propuesta en algo ya empezado.
%
% Material disponible: la campana d303-a6262712 del repositorio (rama feature/fase-estrategias)
% tiene las celdas necesarias. Los numeros se citan por campana/celda/commit, nunca por
% "Decision NNN" (la numeracion 039-051 esta duplicada entre ramas; ver manuscritos/PLAN.md).
\chapter{Experimentos preliminares y resultados}

% PENDIENTE: redactar a partir de la campana d303-a6262712.
% Lo que ya esta verificado contra el ledger y los JSON de la campana:
%   - ancla B0.06: 0,4948 +/- 0,0098, con una mejora del 7,6 % sobre la linea de base;
%   - celda B1.06, tau fijo: empata con el ancla, no la supera;
%   - contraste de Diebold-Mariano contra persistencia: p = 0,16, es decir que con el periodo de
%     evaluacion disponible la diferencia todavia no es distinguible del azar.
% Ese ultimo punto es el mas importante del capitulo y conviene no maquillarlo: motiva el trabajo
% en lugar de debilitarlo, porque muestra que el problema no esta resuelto.

% PENDIENTE: tabla de resultados preliminares, al estilo de la Tab. 3.1 del plan de ejemplo.

% PENDIENTE: enumerar los experimentos siguientes (modulacion de tau, comparacion contra las
% funciones objetivo que devuelva el relevamiento, evaluacion por regimen).

% Plan de tesis -- Capitulo 4.
%
% Correlato en el plan de ejemplo: capitulo 4 (Tiempo estimado de trabajo), media pagina, una sola
% tabla de etapa / tarea / tiempo estimado. No lleva fechas de calendario, solo duraciones.
\chapter{Tiempo estimado de trabajo}

Se estima el tiempo de cada etapa de trabajo sobre la tesis expuesta en el siguiente cuadro.

% PROVISORIO: las duraciones las define el autor. Las etapas siguen el orden de los objetivos
% especificos del capitulo 1.
\begin{center}
\begin{tabular}{@{}lll@{}}
  \toprule
  \textbf{Etapa} & \textbf{Tarea} & \textbf{Tiempo estimado} \\
  \midrule
  Etapa 1 & Relevamiento del estado del arte        & \textit{(falta)} \\
  Etapa 2 & Consolidación del conjunto de datos     & \textit{(falta)} \\
  Etapa 3 & Implementación de las funciones objetivo & \textit{(falta)} \\
  Etapa 4 & Experimentación y contraste estadístico & \textit{(falta)} \\
  Etapa 5 & Escritura y reporte                     & \textit{(falta)} \\
  \bottomrule
\end{tabular}
\end{center}


\bibliografia

\end{document}
